\documentclass[12pt]{article}
\usepackage{amsmath, amssymb, amsthm}
\usepackage{centernot}
\usepackage{geometry}
\geometry{margin=2.5cm}

% Simple manual environments - no amsthm numbering involved
\newenvironment{satz}[1]{%
  \par\vspace{6pt}\noindent\textbf{Theorem #1.}\itshape\space
}{%
  \par\vspace{6pt}
}

\newenvironment{definition}[1]{%
  \par\vspace{6pt}\noindent\textbf{Definition #1.}\itshape\space
}{%
  \par\vspace{6pt}
}

\newenvironment{bemerkung}[1]{%
  \par\vspace{6pt}\noindent\textbf{Remark #1.}\itshape\space
}{%
  \par\vspace{6pt}
}

\renewcommand{\qedsymbol}{$\blacksquare$}
\setcounter{secnumdepth}{0}

\title{7.2 Properties of the Integral}
\date{}

\begin{document}

\maketitle

Convention: $(X,\|\cdot\|)$ normed $\mathbb K$-Banach space, $I=[a,b]$, $a<b$ a compact, perfect
interval.

\begin{satz}{7.2.1}
\begin{enumerate}
    \item Let $(f_n)_{n\in\mathbb N_0}\subset S(I,X)$ be uniformly convergent to $f$. Then $f\in S(I,X)$ and
    \[
    \int_a^b f=\lim_{n\to\infty}\int_a^b f_n.
    \]

    \item Let $(f_n)_{n\in\mathbb N_0}\subset S(I,X)$ and $\sum f_n$ uniformly convergent. Then
    $\displaystyle\sum_{n=0}^{\infty}f_n\in S(I,X)$ and
    \[
    \int_a^b\left(\sum_{n=0}^{\infty}f_n\right)
    =\sum_{n=0}^{\infty}\left(\int_a^b f_n\right).
    \]
\end{enumerate}
\end{satz}

\begin{proof}
\begin{enumerate}
    \item The first part follows from Remark 6.3.9 ($S(I,X)$ is a Banach space). From Remark 7.1.10 it follows:
    \[
    \left\|\int_a^b f-\int_a^b f_n\right\|_X
    \le(b-a)\|f-f_n\|_\infty
    \xrightarrow[n\to\infty]{}0.
    \]

    \item Follows from part~1 applied to the partial sums
    \[
    s_n:=\sum_{k=0}^{n}f_k
    \qquad\text{and}\qquad
    s:=\sum_{k=0}^{\infty}f_k.
    \]
\end{enumerate}
\end{proof}

\begin{bemerkung}{7.2.2}
Theorem 7.2.1 does not hold in general for pointwise convergence. Consider e.g.\ $I=[0,1]$, $X=\mathbb R$ and $f_n\in T(I,X)$ with
\[
f_n(x)=
\begin{cases}
0 & x=0,\\
n & 0<x\le\tfrac{1}{n},\quad n\in\mathbb N,\\
0 & \tfrac{1}{n}<x\le 1.
\end{cases}
\]
Then $f_n\xrightarrow[n\to\infty]{\mathrm{pkt}}0$, but $\int_0^1 f_n=1$ for all $n\in\mathbb N$.
\end{bemerkung}

\begin{satz}{7.2.3}
Let $f\in S(I,X)$. Then $\|f\|_X\in S(I,\mathbb R)$ and
\[
\left\|\int_a^b f\right\|_X\le\int_a^b\|f\|_X\le(b-a)\|f\|_\infty.
\]
\end{satz}

\begin{proof}
Let $(f_n)_{n\in\mathbb N_0}\subset T(I,X)$ be uniformly convergent to $f$. Then $\|f_n\|_X\in T(I,\mathbb R)$ for all $n\in\mathbb N_0$. For every $x\in I$:
\[
\bigl|\|f_n\|_X(x)-\|f\|_X(x)\bigr|
=\bigl|\|f_n(x)\|_X-\|f(x)\|_X\bigr|
\le\|f_n(x)-f(x)\|_X
\le\|f_n-f\|_\infty.
\]
Hence $\|f_n\|_X$ converges uniformly to $\|f\|_X$. Therefore
\[
\|f\|_X\in S(I,\mathbb R).
\]
Let $Z_n:=(x_{n,k})_{k=0}^{m}$ be an associated partition of $I$ for $f_n$ and $M_{n,k}:=(x_{n,k}+x_{n,k-1})/2$ the interval midpoints. Then for every $n\in\mathbb N_0$:
\[
\left\|\int_a^b f_n\right\|_X
=\left\|\sum_{k=1}^{m}f_n(M_{n,k})(x_{k,n}-x_{k-1,n})\right\|_X
\le\sum_{k=1}^{m}\|f_n(M_{n,k})\|_X(x_{k,n}-x_{k-1,n})
=\int_a^b\|f_n\|_X
\le(b-a)\|f_n\|_\infty.\tag{$*$}
\]
From Theorem 7.2.1, passing to the limit $n\to\infty$ on the left-hand side:
\[
\lim_{n\to\infty}\left\|\int_a^b f_n\right\|_X
=\left\|\lim_{n\to\infty}\int_a^b f_n\right\|_X
=\left\|\int_a^b f\right\|_X.
\]
The right-hand side in ($*$) converges to $\|f\|_\infty$, and the claim follows.
\end{proof}

\begin{definition}{7.2.4 (Oriented area)}
Let $a<b$ and $f\in S(I,X)$. Then we define
\[
\int_b^a f:=-\int_a^b f
\qquad\text{and}\qquad
\int_a^a f:=0.
\]
\end{definition}

\begin{satz}{7.2.5}
Let $f\in S(I,X)$ and $a,b,c\in I$. Then:
\[
\int_a^b f=\int_a^c f+\int_c^b f.
\]
\end{satz}

\begin{proof}
Let $a\le b\le c$. Let $(f_n)_{n\in\mathbb N_0}\subset T(I,X)$ be uniformly convergent to $f$. Then
for every compact subinterval $J\subset I$:
\[
T(J,X)\ni f_n|_J\xrightarrow[n\to\infty]{\mathrm{glm}} f_n|_J\in S(J,X).
\]
As one easily verifies, for every $n\in\mathbb N_0$:
\[
\int_a^c f_n=\int_a^b f_n+\int_b^c f_n.
\]
By Theorem 7.2.1 it follows that
\[
\int_a^c f=\int_a^b f+\int_b^c f
\quad\iff\quad
\int_a^b f=\int_a^c f-\int_b^c f=\int_a^c f+\int_c^b f.
\]
\end{proof}

\begin{satz}{7.2.6}
\begin{enumerate}
    \item Let $f\in S(I,\mathbb R)$ and $f\ge 0$. Then $\int_a^b f\ge 0$.

    \item Let $f,g\in S(I,\mathbb R)$ with $f\ge g$. Then $\int_a^b f\ge\int_a^b g$.

    \item Let $f\in S(I,\mathbb R)$ with $f\ge 0$. Suppose there exists $\alpha\in I$ with $f(\alpha)>0$ and $f$ is continuous at $\alpha$.
    Then $\int_a^b f>0$.
\end{enumerate}
\end{satz}

\begin{proof}
\begin{enumerate}
    \item By Remark 6.3.5, there exists a sequence $(f_n)_{n\in\mathbb N_0}\subset T(I,\mathbb R)$ with $f_n\to f$ uniformly
    and $f_n\ge 0$ for all $n\in\mathbb N_0$. Evidently:
    \[
    \int_a^b f_n\ge 0\qquad\forall n\in\mathbb N_0,
    \]
    and the claim for jump-continuous functions follows by passing to the limit $n\to\infty$.

    \item Follows from (1) applied to $f-g$.

    \item By the continuity of $f$ at $\alpha$ and $f(\alpha)>0$, there exists $\delta>0$ with
    \[
    f(x)\ge\tfrac{1}{2}f(\alpha)>0
    \qquad\forall x\in\,]\alpha-\delta,\alpha+\delta[\,\cap I.
    \]
    From Theorem 7.2.5 and parts (1), (2):
    \[
    \int_a^b f
    =\int_a^{\alpha-\delta}f+\int_{\alpha-\delta}^{\alpha+\delta}f+\int_{\alpha+\delta}^{b}f
    \ge\int_{\alpha-\delta}^{\alpha+\delta}\tfrac{1}{2}f(\alpha)
    =\frac{2\delta}{2}f(\alpha)
    =\delta f(\alpha)>0.
    \]
\end{enumerate}
\end{proof}

\begin{satz}{7.2.7 (Integration of vector-valued functions)}
\begin{enumerate}
    \item Let $f:I\to\mathbb K^n$. Then $f\in S(I,\mathbb K^n)$ if and only if for every component $f_i$,
    $1\le i\le n$:
    \[
    f_i\in S(I,\mathbb K).
    \]
    Moreover:
    \[
    \int_a^b f=\left(\int_a^b f_1,\,\int_a^b f_2,\,\ldots,\,\int_a^b f_n\right).
    \]

    \item Let $f=g+ih:I\to\mathbb C$ with $g,h\in I\to\mathbb R$. Then $f\in S(I,\mathbb C)$ if and only if
    $g,h\in S(I,\mathbb R)$. Moreover:
    \[
    \int_a^b f=\int_a^b g+i\int_a^b h.
    \]
\end{enumerate}
\end{satz}

\begin{proof}
\begin{enumerate}
    \item We have:
    \begin{align*}
    f\in S(I,\mathbb K^n)
    &\iff \exists\,(f_m)_{m\in\mathbb N_0}\subset T(I,\mathbb K^n):\ f_m\xrightarrow[m\to\infty]{\mathrm{glm}} f\\
    &\overset{(*)}{\iff} \exists\,(f_m)_{m\in\mathbb N_0}\subset T(I,\mathbb K^n):\ f_{m,i}\xrightarrow[m\to\infty]{\mathrm{glm}} f_i\qquad\forall\,1\le i\le n\\
    &\iff f_i\in S(I,\mathbb K)\qquad\forall\,1\le i\le n.
    \end{align*}
    The statement about the integrals is trivial for step functions. For jump-continuous functions it follows by passing to the limit.

    \item Follows from part~(1) by identifying $\mathbb C$ with $\mathbb R^2$.
\end{enumerate}

\medskip
\noindent\textbf{Ad ($*$):}

\noindent``$\Longrightarrow$''\quad We have:
\begin{align*}
\|f_m-f\|_\infty
&=\sup_{x\in I}\|f_m(x)-f(x)\|_{\mathbb K^n}
\ge C\sup_{x\in I}\max_{1\le i\le n}|f_{m,i}(x)-f_i(x)|\\
&\ge C\sup_{x\in I}|f_{m,i}(x)-f_i(x)|
=C\|f_{m,i}-f_i\|_\infty.
\end{align*}

\noindent``$\Longleftarrow$''\quad
\[
\forall\varepsilon>0\ \exists\,M_i>0\ \forall x\in I\ \forall m\ge M_i
\implies |f_{i,m}(x)-f_i(x)|\le\varepsilon,
\qquad\forall\,1\le i\le n.
\]
Set $M:=\max_{1\le i\le n}M_i$. Then:
\[
\forall\varepsilon>0\ \forall m\ge M
\implies \varepsilon\ge\sup_{x\in I}\max_{1\le i\le n}|f_{i,m}(x)-f_i(x)|
=\sup_{x\in I}\|f_m(x)-f(x)\|_{\mathbb K^n}
\ge c\,\|f_m-f\|_\infty,
\]
i.e.\ $f_m\xrightarrow[m\to\infty]{\mathrm{glm}} f$.
\end{proof}

\begin{satz}{7.2.8 (Dependence of the integral on the limits of integration)}
Let $f\in S(I,X)$ and $c\in I$. Then the function
\[
F(x):=\int_c^x f\qquad\forall x\in I
\]
belongs to $C(I,X)$ and satisfies:
\[
\|F(x)-F(y)\|_X\le|x-y|\,\|f\|_\infty
\qquad\forall x,y\in I.
\]
\end{satz}

\begin{proof}
From Theorem 7.2.5 it follows for $x,y\in I$:
\[
F(x)-F(y)=\int_c^x f-\int_c^y f=\int_y^x f
\]
and with Remark 7.1.10:
\[
\|F(x)-F(y)\|_X\le\|f\|_\infty|y-x|.
\]
\end{proof}

\begin{satz}{7.2.9}
Let $f\in S(I,X)$ and continuous at $\alpha\in I$. Then the function
\[
F(x):=\int_a^x f\qquad\forall x\in I
\]
is differentiable at $\alpha$, and:
\[
F'(\alpha)=f(\alpha).
\]
\end{satz}

\begin{proof}
Theorem 7.2.8 implies: $F$ is continuous. Let $\varepsilon>0$ be arbitrary. Then there exists $\delta>0$ with:
\[
\|f(\xi)-f(\alpha)\|_X<\varepsilon\qquad\forall\xi\in I:\ |\xi-\alpha|\le\delta.
\]
Let $h\in\mathbb R\setminus\{0\}$ with $\alpha+h\in I$ and $|h|<\delta$. Then:
\begin{align*}
\left\|\frac{F(\alpha+h)-F(\alpha)}{h}-f(\alpha)\right\|_X
&=\left\|\left(\frac{1}{h}\int_\alpha^{\alpha+h}f\right)-f(\alpha)\right\|_X
=\left\|\left(\frac{1}{h}\int_\alpha^{\alpha+h}f\right)-\frac{1}{h}\int_\alpha^{\alpha+h}f(\alpha)\right\|_X\\
&=\frac{1}{|h|}\left\|\int_\alpha^{\alpha+h}(f-f(\alpha))\right\|_X
\le\frac{1}{|h|}\left\|\int_\alpha^{\alpha+h}\varepsilon\right\|_X
=\frac{1}{|h|}\cdot|h|\cdot\varepsilon=\varepsilon.
\end{align*}
\end{proof}

\begin{definition}{7.2.10}
Let $f,F:I\to X$ be two functions. $F$ is called an \emph{antiderivative} of $f$ if
$F$ is differentiable at every point of $I$ and
\[
F'(x)=f(x)\qquad\forall x\in I.
\]
The set of all antiderivatives of $f$ is denoted by $\int f$ or $\int f(x)\,dx$ and is called the
\emph{indefinite integral} of $f$.
\end{definition}

\begin{bemerkung}{7.2.11}
Let $F,G$ be two antiderivatives of $f$. Then $F-G$ is constant on $I$.
\end{bemerkung}

\begin{proof}
Follows from Theorem 5.2.6 ($(F-G)'=F'-G'=f-f=0\implies F-G$ constant). \qedhere
\end{proof}

\begin{satz}{7.2.12 (Fundamental theorem of calculus)}
Every $f\in C(I,X)$ possesses an antiderivative $F$ of $f$ and satisfies:
\[
F(x)=F(a)+\int_a^x f=F(a)+\int_a^x F'
\qquad\forall x\in I.
\]
In particular, for every antiderivative $F$ of $f$:
\[
\int_a^b f=F(b)-F(a)=:F\big|_a^b.
\]
\end{satz}

\begin{proof}
Let $G(x):=\int_a^x f$. By Theorem 7.2.9, $G$ is an antiderivative of $f$. From
Remark 7.2.11 it follows for every antiderivative $F$ of $f$:
\begin{align*}
F(x)&=G(x)+c,\qquad c\in X,\ \forall x\in I,\\
F(a)&=G(a)+c=c
\end{align*}
\[
\implies F(x)=F(a)+\int_a^x f\qquad\forall x\in I.
\]
\end{proof}

\par\vspace{6pt}\noindent\textbf{Example 7.2.13.}\itshape\space
\begin{enumerate}
    \item Table of antiderivatives:

    \upshape
    \[
    \renewcommand{\arraystretch}{1.4}
    \begin{array}{c|c|cc}
    f & \int f & & \\\hline
    x^\alpha & \frac{1}{\alpha+1}x^{\alpha+1} & +c & \alpha\neq-1\\
    \frac{1}{x} & \ln|x| & +c & \\
    a^x & \frac{1}{\ln a}a^x & +c & a>0\\
    e^{ax} & \frac{1}{a}e^{ax} & +c & a\in\mathbb C\setminus\{0\}\\
    \cos(x) & \sin(x) & +c & \\
    -\sin(x) & \cos(x) & +c & \\
    \frac{1}{\cos^2 x} & \tan(x) & +c & \\
    \frac{-1}{\sin^2 x} & \cot(x) & +c & \\
    \frac{1}{\sqrt{1-x^2}} & \arcsin x & +c & \\
    \frac{1}{1+x^2} & \arctan x & +c & \\
    \sum_{k=0}^{\infty}a_k(x-x_0)^k & \sum_{k=0}^{\infty}\frac{a_k}{k+1}(x-x_0)^{k+1} & +c &
    \end{array}
    \]
    \itshape

    \item $f\in C^1(I,\mathbb R)$, $f(x)\neq 0$ for all $x\in I$
    \[
    \implies\int\frac{f'}{f}=\ln(|f|)+c.
    \]
\end{enumerate}
\upshape

\begin{proof}
\begin{enumerate}
    \item Follows by differentiation of the right-hand side.

    \item By the intermediate value theorem 4.4.7, for all $x\in I$ either
    \begin{enumerate}
        \item[(a)] $f(x)>0$, or
        \item[(b)] $f(x)<0$.
    \end{enumerate}

    \textbf{Case (a):} $(\ln(|f|))'=(\ln f)'=\dfrac{f'}{f}$.

    \textbf{Case (b):} $(\ln(|f|))'=(\ln(-f))'=\dfrac{-f'}{-f}=\dfrac{f'}{f}$.
\end{enumerate}
\end{proof}

\noindent Analogue of the mean value theorem 5.2.5:

\begin{satz}{7.2.14 (Mean value theorem of integral calculus)}
Let $f,\phi\in C(I,\mathbb R)$ with $\phi\ge 0$ on $I$. Then there exists $\xi$, $a\le\xi\le b$:
\[
\int_a^b f\cdot\phi=f(\xi)\int_a^b\phi.
\]
In particular for $\phi\equiv 1$: there exists $\eta$, $a\le\eta\le b$:
\[
\bigl(F(b)-F(a)=\bigr)\qquad\int_a^b f=f(\eta)(b-a)\qquad\bigl(=F'(\eta)(b-a)\bigr).
\]
\end{satz}

\begin{proof}
If $\int_a^b\phi=0$, then it is an exercise to show that $\phi\equiv 0$. In this case
the claim holds trivially. So let $\int_a^b\phi>0$. Define
\[
m:=\inf_{x\in I}f(x)\qquad M:=\sup_{x\in I}f(x).
\]
Then
\[
m\phi(x)\le f(x)\phi(x)\le M\phi(x)\qquad\forall x\in I
\]
and by Theorem 7.2.6:
\[
m\int_a^b\phi\le\int_a^b f\phi\le M\int_a^b\phi
\qquad\implies\qquad
m\le\frac{\int_a^b f\phi}{\int_a^b\phi}\le M.
\]
The intermediate value theorem 4.4.7 implies: there exists $\xi\in\,]a,b[\,$ with
\[
f(\xi)=\frac{\int_a^b f\cdot\phi}{\int_a^b\phi}.
\]
This proves the first claim. The second claim follows from the first with the special choice $\phi=1$.
\end{proof}

\end{document}